\documentclass[a4paper,10pt]{article}
\usepackage[utf8]{inputenc}
\usepackage[T1]{fontenc}
\usepackage[french]{babel}
\usepackage[margin=1.5cm]{geometry}
\usepackage{xcolor}
\usepackage{tcolorbox}
\usepackage{enumitem}
\usepackage{multicol}
\usepackage{lmodern}
\usepackage{tabularx}
\usepackage{booktabs}
\usepackage{amsmath}

\renewcommand{\familydefault}{\sfdefault}
\pagestyle{empty}
\setlist{nosep, leftmargin=1.5em}

\definecolor{navy}{RGB}{0,51,102}
\definecolor{teal}{RGB}{0,120,120}
\definecolor{crimson}{RGB}{180,0,0}
\definecolor{green}{RGB}{0,120,50}
\definecolor{lightbg}{RGB}{245,248,252}
\definecolor{orange}{RGB}{200,100,0}

\tcbset{basebox/.style={
  colback=lightbg, colframe=navy, coltitle=white,
  fonttitle=\bfseries, arc=2mm, boxrule=0.5mm,
  left=3mm, right=3mm, top=2mm, bottom=2mm, after skip=6pt
}}

\newcommand{\key}[1]{\textbf{\textcolor{navy}{#1}}}
\newcommand{\warn}[1]{\textbf{\textcolor{crimson}{#1}}}
\newcommand{\tip}[1]{\textit{\textcolor{teal}{#1}}}

\begin{document}

\begin{center}
  \colorbox{navy}{\parbox{\dimexpr\textwidth-2\fboxsep}{
    \centering\vspace{3pt}
    {\LARGE\bfseries\textcolor{white}{RÉSEAUX \& ROUTAGE — Fiche de révision détaillée}}\\
    {\normalsize\textcolor{white}{BTS SIO — Maël Mignard}}
    \vspace{3pt}
  }}
\end{center}
\vspace{6pt}

\begin{multicols}{2}

% ============================================================
\begin{tcolorbox}[basebox, colframe=navy, title=1. Le Modèle OSI — 7 couches]
{\footnotesize
\begin{tabularx}{\columnwidth}{@{}clXl@{}}
\toprule
\textbf{N°} & \textbf{Couche} & \textbf{Rôle} & \textbf{PDU} \\
\midrule
7 & Appli. & Interface utilisateur & Message \\
6 & Présent. & Format, chiffrement & Message \\
5 & Session & Gestion des sessions & Message \\
4 & Transport & Fiabilité (TCP/UDP) & Segment \\
3 & Réseau & Adressage/routage IP & Paquet \\
2 & Liaison & Trames, MAC, switch & Trame \\
1 & Physique & Signaux élec./opt. & Bits \\
\bottomrule
\end{tabularx}}
\vspace{4pt}\\
\key{Moyen mnémotechnique (fr)} : « \textit{Ah Petit Salaud Tu Rends La Peine !} » \\
\key{Encapsulation} : en émission, chaque couche ajoute son en-tête. \\
\key{Désencapsulation} : en réception, chaque couche retire son en-tête.
\end{tcolorbox}

% ============================================================
\begin{tcolorbox}[basebox, colframe=teal, title=2. Modèle TCP/IP (4 couches)]
{\footnotesize
\begin{tabularx}{\columnwidth}{@{}lX@{}}
\toprule
\textbf{Couche TCP/IP} & \textbf{Équivalent OSI} \\
\midrule
Application  & OSI 5, 6, 7 \\
Transport    & OSI 4 (TCP/UDP) \\
Internet     & OSI 3 (IP, ICMP) \\
Accès réseau & OSI 1, 2 \\
\bottomrule
\end{tabularx}}
\vspace{4pt}\\
\tip{Différence OSI/TCP-IP} : OSI est un modèle \textbf{théorique} de référence, TCP/IP est le \textbf{modèle pratique} utilisé sur Internet.
\end{tcolorbox}

% ============================================================
\begin{tcolorbox}[basebox, colframe=orange, title=3. TCP vs UDP]
\begin{tabularx}{\columnwidth}{lXX}
\toprule
 & \textbf{TCP} & \textbf{UDP} \\
\midrule
Connexion & Oui (3-way handshake) & Non \\
Fiabilité & Garantie (ACK) & Aucune \\
Vitesse & Plus lent & Rapide \\
Ordre & Garanti & Non garanti \\
Usage & HTTP, FTP, SSH & DNS, VoIP, streaming \\
\bottomrule
\end{tabularx}
\vspace{4pt}\\
\key{3-Way Handshake TCP} : SYN $\to$ SYN-ACK $\to$ ACK
\end{tcolorbox}

% ============================================================
\begin{tcolorbox}[basebox, colframe=green, title=4. Adressage IPv4]
\key{Format} : 4 octets (32 bits) — ex: \texttt{192.168.1.100/24}\\[3pt]
\textbf{Plages privées (RFC 1918) :}
\begin{itemize}
  \item \texttt{10.0.0.0/8} \hfill (Classe A — grandes entreprises)
  \item \texttt{172.16.0.0/12} \hfill (Classe B)
  \item \texttt{192.168.0.0/16} \hfill (Classe C — particuliers)
  \item \texttt{127.0.0.1} \hfill (Loopback)
\end{itemize}

\textbf{Masques et taille de réseau :}\\
{\footnotesize
\begin{tabularx}{\columnwidth}{@{}llX@{}}
/24 & 255.255.255.0   & 254 hôtes \\
/25 & 255.255.255.128 & 126 hôtes \\
/16 & 255.255.0.0     & 65\,534 hôtes \\
/8  & 255.0.0.0       & 16M hôtes \\
\end{tabularx}}\\[3pt]
\key{Calcul hôtes} : $2^{(32 - \text{préfixe})} - 2$ \\
\warn{-2} : adresse réseau + broadcast \\[3pt]
\key{NAT} (Network Address Translation) : traduit IP privée $\to$ publique \\
\key{DHCP} : attribue automatiquement IP, masque, passerelle, DNS
\end{tcolorbox}

% ============================================================
\begin{tcolorbox}[basebox, colframe=green, title=5. IPv6 — Essentiel]
\key{Format} : 128 bits, 8 groupes hexadécimaux \\
Ex: \texttt{2001:0db8:85a3::8a2e:0370:7334} \\[3pt]
\textbf{Types d'adresses :}
\begin{itemize}
  \item \key{Unicast} : une seule destination
  \item \key{Multicast} : groupe de destinations
  \item \key{Anycast} : le plus proche du groupe
  \item \key{Lien-local} : \texttt{fe80::/10} (non routable)
\end{itemize}
\textbf{Avantages sur IPv4} : espace d'adressage immense, pas de NAT nécessaire, sécurité intégrée (IPsec).
\end{tcolorbox}

% ============================================================
\begin{tcolorbox}[basebox, colframe=navy, title=6. Protocoles applicatifs — Ports]
{\footnotesize
\begin{tabularx}{\columnwidth}{@{}llX@{}}
\toprule
\textbf{Proto} & \textbf{Port} & \textbf{Usage} \\
\midrule
HTTP  & 80    & Web non chiffré \\
HTTPS & 443   & Web chiffré (TLS) \\
FTP   & 21    & Transfert fichiers \\
SSH   & 22    & Shell sécurisé \\
SMTP  & 25    & Envoi mail \\
DNS   & 53    & Résolution noms \\
DHCP  & 67/68 & Attribution IP \\
RDP   & 3389  & Bureau à distance \\
SNMP  & 161   & Supervision réseau \\
\bottomrule
\end{tabularx}}
\end{tcolorbox}

% ============================================================
\begin{tcolorbox}[basebox, colframe=teal, title=7. Équipements réseau]
{\footnotesize
\begin{tabularx}{\columnwidth}{@{}llX@{}}
\toprule
\textbf{Équipement} & \textbf{L.} & \textbf{Rôle} \\
\midrule
Hub           & 1   & Diffuse tout (obsolète) \\
Switch        & 2   & Commute par adresse MAC \\
Routeur       & 3   & Route par adresse IP \\
Firewall      & 3-7 & Filtre le trafic \\
Point d'accès & 1-2 & Connexion WiFi \\
\bottomrule
\end{tabularx}}
\vspace{4pt}\\
\key{Table MAC} : switch mémorise MAC $\leftrightarrow$ port \\
\key{Table ARP} : résolution IP $\leftrightarrow$ MAC \\
\texttt{arp -a} : afficher la table ARP
\end{tcolorbox}

% ============================================================
\begin{tcolorbox}[basebox, colframe=orange, title=8. VLAN — Virtual LAN]
\key{VLAN} : segmentation \textbf{logique} d'un réseau physique \\[3pt]
\textbf{Avantages :}
\begin{itemize}
  \item Isolation du trafic (sécurité)
  \item Réduction du domaine de broadcast
  \item Flexibilité sans recâblage
\end{itemize}
\vspace{3pt}
\key{VLAN d'accès} : port connecté à un terminal (1 VLAN) \\
\key{Lien trunk (802.1Q)} : port inter-switch (plusieurs VLAN) \\[3pt]
\textbf{Configuration Cisco :}
\small
\begin{verbatim}
Switch(config)# vlan 10
Switch(config-vlan)# name IT
Switch(config)# interface Fa0/1
Switch(config-if)# switchport mode access
Switch(config-if)# switchport access vlan 10
\end{verbatim}
\normalsize
\end{tcolorbox}

% ============================================================
\begin{tcolorbox}[basebox, colframe=crimson, title=9. Spanning Tree Protocol (STP)]
\key{Problème résolu} : éviter les boucles de commutation \\[3pt]
\textbf{Fonctionnement :}
\begin{itemize}
  \item Élection d'un \key{Root Bridge} (plus faible Bridge ID)
  \item Ports en état \textit{Forwarding} ou \textit{Blocking}
  \item Si un lien tombe, recalcul automatique
\end{itemize}
\vspace{3pt}
\key{RSTP (802.1w)} : version rapide du STP \\
\key{PVST+} : une instance STP par VLAN (Cisco)
\end{tcolorbox}

% ============================================================
\begin{tcolorbox}[basebox, colframe=green, title=10. Routage — Concepts]
\key{Routage statique} : routes saisies manuellement \\
\quad $\to$ Simple, adapté aux petits réseaux \\
\key{Routage dynamique} : protocoles qui échangent les routes \\[3pt]
\textbf{Protocoles de routage :}
\begin{itemize}
  \item \key{RIP} : distance en sauts (max 15), lent
  \item \key{OSPF} : coût (bande passante), rapide
  \item \key{BGP} : Internet, entre AS (opérateurs)
\end{itemize}
\vspace{3pt}
\textbf{Commandes Cisco :}
\begin{verbatim}
Router(config)# ip route 192.168.2.0
              255.255.255.0 10.0.0.1
Router# show ip route
\end{verbatim}
\end{tcolorbox}

% ============================================================
\begin{tcolorbox}[basebox, colframe=navy, title=11. Multicast]
\key{Unicast} : 1 émetteur $\to$ 1 récepteur \\
\key{Broadcast} : 1 émetteur $\to$ tous \\
\key{Multicast} : 1 émetteur $\to$ groupe abonné \\[3pt]
\textbf{Plage multicast IPv4} : \texttt{224.0.0.0/4} \\
\key{IGMP} : protocole de gestion des groupes multicast \\
\key{Usage} : streaming vidéo, IPTV, protocoles de routage
\end{tcolorbox}

% ============================================================
\begin{tcolorbox}[basebox, colframe=teal, title=12. Commandes de diagnostic réseau]
{\footnotesize
\begin{tabularx}{\columnwidth}{@{}lX@{}}
\toprule
\textbf{Commande} & \textbf{Usage} \\
\midrule
\texttt{ping 8.8.8.8}          & Test connectivité ICMP \\
\texttt{traceroute/tracert}     & Chemin des paquets \\
\texttt{ipconfig / ip addr}     & Config. IP locale \\
\texttt{netstat -an}            & Connexions actives \\
\texttt{nslookup domaine}       & Résolution DNS \\
\texttt{show interfaces}        & Cisco : état interfaces \\
\texttt{show vlan brief}        & Cisco : liste VLANs \\
\texttt{show mac-address-table} & Cisco : table MAC \\
\bottomrule
\end{tabularx}}
\end{tcolorbox}

\end{multicols}

\vspace{4pt}
\begin{tcolorbox}[basebox, colframe=crimson, title=Points clés à retenir pour l'examen]
\begin{multicols}{3}
\begin{itemize}
  \item Le modèle OSI a \textbf{7 couches}, TCP/IP en a \textbf{4}
  \item TCP = fiable (ACK), UDP = rapide (pas d'ACK)
  \item VLAN = segmentation logique, trunk = lien multi-VLAN
  \item STP évite les boucles en bloquant des ports
  \item /24 = 254 hôtes (255.255.255.0)
  \item Adresses privées : 10.x, 172.16-31.x, 192.168.x
  \item OSPF > RIP (plus rapide, pas de limite 15 sauts)
  \item Un routeur travaille en couche 3, un switch en couche 2
  \item MAC = couche 2, IP = couche 3
  \item DNS port 53, HTTPS port 443, SSH port 22
  \item Ping utilise ICMP (couche 3)
  \item NAT traduit IP privée/publique
\end{itemize}
\end{multicols}
\end{tcolorbox}

\end{document}
